
\documentclass[11pt,a4paper]{article}
\usepackage[slovak]{babel}
\usepackage[utf8]{inputenc}
\usepackage[T1]{fontenc}
 
\title{Spresnenie zamerania projektu}
\author{
Dmytro Kolodeznyi \\
[2pt]
{\small Slovenská technická univerzita v Bratislave} \\
{\small Fakulta informatiky a informačných technológií} \\[4pt]
{\small Cvičiaci: Ing. Oliver Udvardi}
}
\date{3. októbra 2026}
 
\begin{document}
\maketitle
 
\section*{Názov projektu}
Detekcia prania špinavých peňazí v kryptomenových transakciách
 
\section*{Akronym}
AMLGRAPH (Analysis of Money Laundering via Graph-based Recognition of Anomalous Patterns in bHlockchain)
 
\section*{Anotácia}
Blockchainové siete sú verejné a nemenné, preto je každá transakcia dohľadateľná, no identita vlastníkov adries zostáva skrytá. Túto medzeru páchatelia zneužívajú na pranie špinavých peňazí prostredníctvom mixerov, dlhých reťazcov transakcií a zmenární bez overenia totožnosti.
 
Projekt AMLGRAPH sa zameriava na analýzu štruktúry transakčného grafu kryptomenových sietí s cieľom odhaliť typické vzorce prania špinavých peňazí. Preskúmame heuristické metódy (fan-in/fan-out vzorce, peel chains) aj metódy strojového učenia nad grafmi (node embedding, grafové neurónové siete), vrátane dostupných datasetov (napr. Elliptic).
 
Téma má silnú normatívnu oporu v nariadení MiCA, smerniciach o boji proti praniu špinavých peňazí a odporúčaniach FATF (Travel Rule). Výsledky budú merateľné pomocou metrík precision, recall, miery falošne pozitívnych výsledkov a objemu identifikovaných prostriedkov.
 
Súčasťou práce bude aj diskusia o rovnováhe medzi súkromím používateľov a potrebou regulačného dohľadu nad finančnými tokmi v decentralizovanom prostredí.
 
\end{document}
 